\begin{figure}[tb]
\centering
\begin{tikzpicture}[x=1.45cm,y=1.15cm]
  \draw[gray!45] (-.55,0)--(2.55,0) (-.55,2)--(2.55,2);
  \draw[gray!45] (0,-.5)--(0,2.5) (2,-.5)--(2,2.5);
  \node[draw,circle,fill=white,inner sep=2.2pt] (y1) at (0,0) {};
  \node[draw,circle,fill=white,inner sep=2.2pt] (y2) at (2,2) {};
  \node[circle,fill=blue!65!black,inner sep=2.5pt] (x1) at (2,0) {};
  \node[circle,fill=blue!65!black,inner sep=2.5pt] (x2) at (0,2) {};
  \node[below left=3pt of y1] {\(y_1=(a,b)\)};
  \node[above right=3pt of y2] {\(y_2=(a',b')\)};
  \node[below right=3pt of x1] {\((a,b')\)};
  \node[above left=3pt of x2] {\((a',b)\)};
  \node[rotate=90,fill=white,inner sep=2pt] at (0,1) {\(Y=0\)};
  \node[rotate=90,fill=white,inner sep=2pt] at (2,1) {\(Y=0\)};
  \node[fill=white,inner sep=2pt] at (1,0) {\(Y=1\)};
  \node[fill=white,inner sep=2pt] at (1,2) {\(Y=1\)};
\end{tikzpicture}
\caption{Two candidate center states differing in both coordinates admit
exactly two point states in their common strict cross. The filled points are
those opposite corners. Within this four-state configuration, the bit of a
consistent edge and the state at either endpoint determine the state at
the other endpoint. Rows have fixed first label and columns have fixed
second label. The diagram displays coordinate
equalities of states, not geometric positions of the original plane points.}
\label{com:fig-crosses}
\end{figure}
